\documentclass[]{aastex701}
\usepackage[style=numeric-comp, autocite=superscript, backend=biber]{biblatex}
\usepackage{booktabs}
\usepackage{pifont}
\newcommand{\cmark}{\ding{51}}
\addbibresource{PHYD72_Reeves.bib}
\AtBeginDocument{\setcitestyle{super,numbers,sort&compress}}


\AtBeginDocument{\renewcommand{\bibnumfmt}[1]{[#1]\hspace{6pt}}}

\setlength{\parskip}{1em}

\newcommand{\vdag}{(v)^\dagger}
\newcommand\aastex{AAS\TeX}
\newcommand\latex{La\TeX}
\renewcommand{\sectionautorefname}{Section}
\renewcommand{\subsectionautorefname}{Subsection}
\renewcommand{\subsubsectionautorefname}{Subsubsection}
\begin{document}

\title{Artificial Intelligence in Physics Education: Opportunities for Students with Disabilities}

\author{Taylor Reeves}
\affiliation{University of Toronto}
\email{taylor.reeves@mail.utoronto.ca}

\begin{abstract}
Students with disabilities face substantial barriers in undergraduate physics education, where traditional instructional design often assumes standardized modes of learning that do not accommodate diverse cognitive profiles. Structural features of physics curricula---including high-stakes assessments, dense technical readings, and rapid conceptual progression---can disproportionately disadvantage students with learning disabilities such as dyslexia, dyscalculia, ADHD, and Autism Spectrum Disorder. Recent advances in artificial intelligence (AI) offer promising tools to address these inequities. This paper examines the intersection of disability, physics education, and AI in higher education, drawing on recent literature in physics education research, special education, and educational technology. We analyze how structural, cognitive, and social barriers manifest for disabled students in physics programs, survey current AI applications used by both instructors and students, and evaluate how AI tools---including adaptive tutoring systems, multimodal visualizations, and large language models---can reduce cognitive load, support reading comprehension, and provide flexible, personalized learning environments. Particular attention is given to how AI can scaffold complex multi-step problem solving and make abstract concepts such as quantum mechanics more accessible through dynamic representation. We also examine key concerns, including overreliance, algorithmic bias, AI-generated misinformation, and the disproportionate impact of automated plagiarism detection on students who use assistive technologies. The paper concludes with five evidence-based recommendations for the responsible and inclusive integration of AI into undergraduate physics education, emphasizing structured implementation, instructor training, institutional policy, and ongoing evaluation for equity and accessibility.
\end{abstract}

\keywords{Physics Education --- Artificial Intelligence --- Accessibility --- Learning Disabilities}

\section{Introduction}

Physics occupies a central position within Science, Technology, Engineering, and Mathematics (STEM) education and is widely regarded as one of the most conceptually and mathematically demanding disciplines. Learning physics requires students to integrate multiple forms of reasoning simultaneously, including mathematical modeling, conceptual interpretation of physical systems, and translating real-world phenomena into abstract theoretical frameworks. Unlike many other subjects, physics often demands that students interpret textual descriptions of physical systems, convert these descriptions into mathematical equations, and apply quantitative reasoning to evaluate solutions. This combination of conceptual, mathematical, and linguistic demands can present substantial challenges for students whose learning processes differ from the instructional assumptions embedded in traditional physics courses \autocite{Chasen2025, Gull2025}.

For students with disabilities, particularly those with learning differences affecting reading comprehension, information processing, or executive functioning, the structure of many physics courses can create barriers that extend beyond the inherent complexity of the discipline. Research examining accessibility in physics education indicates that many courses are designed around standardized modes of learning, assuming rapid integration of technical reading, mathematical reasoning, and conceptual abstraction \autocite{Chasen2025, AlAzawei2023}. Physics curricula often emphasize dense textbooks, high-stakes problem-solving assessments, and fast progression through complex material. These instructional practices may disadvantage students who require alternative learning strategies, additional processing time, or diverse forms of explanation. Consequently, observed differences in academic outcomes between disabled and non-disabled students may reflect structural features of educational environments rather than differences in intellectual ability or scientific potential.

Recent advances in artificial intelligence (AI) technologies have introduced new possibilities for addressing some of these challenges within physics education. AI systems are increasingly capable of generating explanations, visualizing complex physical phenomena, and providing interactive tutoring for problem solving. These tools can present scientific concepts through multiple representations, including diagrams, simulations, and step-by-step derivations. Such multimodal learning approaches are particularly valuable in physics, where abstract concepts---such as electromagnetic fields, orbital motion, or quantum wavefunctions---are difficult to interpret through equations alone. Research in physics education and special education indicates that interactive and visual representations can improve conceptual understanding, particularly for students who benefit from spatial or multimodal learning strategies \autocite{Voultsiou2025, HarkinsBrown2025, Singh2001}.

In addition to supporting conceptual visualization, AI technologies can assist students throughout multi-stage problem-solving processes. Physics problems frequently require identifying relevant principles, translating textual descriptions into mathematical expressions, manipulating equations, and interpreting the physical significance of results. For students with learning disabilities, these tasks can impose a substantial cognitive burden. AI-based tutoring systems can reduce this load by breaking problems into smaller components, offering targeted hints, and providing alternative explanations when students encounter difficulties. When implemented thoughtfully, such systems may help students build deeper conceptual understanding while maintaining engagement with challenging material \autocite{Paglialunga2025, Sweller2011}.

However, the growing integration of AI into physics education also raises important questions regarding its appropriate role. While AI can enhance accessibility and conceptual understanding, it introduces potential risks related to overreliance, misinformation, algorithmic bias, and academic integrity. These concerns may be amplified for students with disabilities, who may depend more heavily on assistive technologies and therefore experience both the advantages and limitations of AI more acutely than other learners. Careful evaluation and responsible implementation are necessary to ensure that AI supports learning without undermining the development of independent reasoning skills.

This paper explores the intersection of disability, physics education, and artificial intelligence in higher education. We examine how structural, cognitive, and social factors shape the experiences of students with disabilities in physics courses in \autoref{sec:disability_STEMEd} and identify current AI applications for instructors and students in \autoref{sec:AI_STEMEd}. We analyze how AI tools can support learning through adaptive tutoring, visualization, and flexible problem-solving strategies in \autoref{sec:AI_Support}. In \autoref{sec:concerns} we highlight potential limitations and risks associated with AI integration, but in \autoref{sec:discussion} we discuss and summarize how these technologies may contribute to more accessible and inclusive physics learning environments while highlighting the challenges that accompany their adoption. We provide some recommendations for instructors and learners in \autoref{sec:recommendations}.



\section{How Do Disabilities Manifest in Physics Education}
\label{sec:disability_STEMEd}
Students with disabilities encounter a variety of structural, cognitive, and social barriers in higher education, particularly within physics programs. Refs. \cite{Chasen2025} and \cite{Gull2025} note that although universities have increasingly adopted accessibility services and accommodations, many physics courses continue to assume a standardized mode of learning that does not reflect the diversity of student needs. Ref.~\cite{Chasen2025} further emphasizes that students with disabilities are often evaluated against norms developed for non-disabled peers, rather than within a framework that recognizes differences in learning processes and cognitive styles. This comparison can unintentionally reinforce inequities and contribute to negative educational experiences in STEM disciplines.

A frequently highlighted challenge is the persistence of comparative evaluation practices within physics education. Students with disabilities are commonly assessed using traditional academic expectations designed for non-disabled learners \autocite{Chasen2025, AlAzawei2023}. Ref.~\cite{Chasen2025} found that the rigid structure of physics curricula, including high-stakes examinations, dense technical readings, and rapid progression through concepts, can exacerbate accessibility barriers. Ref.~\cite{Almufareh2025} similarly notes that such structural features may disadvantage students who require alternative learning strategies, extended processing time, or modified instructional approaches. Consequently, lower academic performance among disabled students often reflects misaligned instructional design rather than differences in ability.

Reading comprehension is another significant barrier, particularly for students with learning disabilities. Paglialunga and Melogno \autocite{Paglialunga2025} and Voultsiou and Moussiades \autocite{Voultsiou2025} highlight that physics education relies heavily on technical textbooks and dense textual explanations. Students with Attention Deficit-Hyperactivity Disorder (ADHD), Autism Spectrum Disorder (ASD), dyslexia, or dyscalculia may struggle with sustained attention or decoding complex passages \autocite{Paglialunga2025, Voultsiou2025, HarkinsBrown2025}. Kontopoulou, Drigas, and Reisis \autocite{Kontopoulou2024} emphasize that these difficulties can affect students' ability to interpret abstract terminology or maintain focus on extended readings, which is critical in physics where conceptual understanding often depends on precise interpretation of both text and mathematical expressions. As a result, students may face substantial obstacles in building foundational knowledge \autocite{Paglialunga2025, Sweller2011}.

The challenges of reading-related disabilities are particularly pronounced when students must simultaneously process linguistic and mathematical information. Dyslexia and dyscalculia can interfere with the integration of textual problem descriptions and quantitative reasoning \autocite{Kontopoulou2024}. Ref.~\cite{Kaliampos2025} describes how students with dyslexia may struggle with reading components of physics problems, while students with dyscalculia may face difficulties manipulating numerical relationships or mathematical symbols. These overlapping challenges can compound when courses assume rapid integration of reading comprehension and mathematical reasoning, further hindering learning outcomes \autocite{Kontopoulou2024}. Table~\ref{tab:disability_barriers} summarizes how specific learning disabilities interact with core physics tasks.

\begin{table}[h!]
\centering
\renewcommand{\arraystretch}{1.4}
\setlength{\tabcolsep}{6pt}
\resizebox{\textwidth}{!}{
\begin{tabular}{lcccccc}
\toprule
\textbf{Disability}
& \textbf{Reading}
& \textbf{Math}
& \textbf{Attention}
& \textbf{Abstract Concepts}
& \textbf{Multi-step Problems}
& \textbf{Social / Self-Advocacy} \\
\midrule
Dyslexia
& \cmark &  & \cmark & \cmark & \cmark & \cmark \\
Dyscalculia
&  & \cmark &  & \cmark & \cmark &  \\
ADHD
& \cmark &  & \cmark &  & \cmark & \cmark \\
ASD
& \cmark &  & \cmark & \cmark &  & \cmark \\
\bottomrule
\end{tabular}
}
\caption{Summary of common challenges faced by students with different disabilities in physics education.
\textbf{Note:} \cmark\ indicates an area where students commonly experience difficulty. Adapted from \autocite{Kontopoulou2024, Kaliampos2025, Paglialunga2025}.}
\label{tab:disability_barriers}
\end{table}

Continuous self-advocacy represents another barrier for students with disabilities, and its demands differ markedly between secondary and tertiary education. In secondary school, disability accommodations are often coordinated by institutional support staff, with parents or guardians playing an active role in ensuring students receive appropriate services. In contrast, post-secondary institutions place the full responsibility of self-identification and accommodation-seeking on the student \autocite{Gull2025, AlAzawei2023}. While disability services provide essential accommodations at the university level, students must actively request extended time, alternative formats, or flexible deadlines --- often on a course-by-course basis. This process can be emotionally and cognitively taxing, particularly in physics programs where instructors may lack training in adapting teaching methods for neurodiverse learners \autocite{Almufareh2025, Kaliampos2025}. As a result, students may repeatedly need to explain their circumstances in academic contexts that are not fully equipped to support diverse learning needs, a burden that compounds across a degree program.

Social stigma further influences educational experiences. Disabilities are often invisible, and students may fear negative judgments from peers or instructors if they disclose their learning needs \autocite{Gull2025, AlAzawei2023}. This concern can reduce engagement in collaborative problem solving or classroom discussions, as students may perceive requesting support as a sign of weakness \autocite{Gull2025, AlAzawei2023}. Such social dynamics contribute to inequities in participation and learning outcomes \autocite{AlAzawei2023}.

The pace and structure of STEM courses also exacerbate accessibility challenges. Chasen et al.\ \autocite{Chasen2025} and Almufareh et al.\ \autocite{Almufareh2025} note that rapid progression through complex material leaves limited opportunities for reviewing concepts or practicing foundational skills. Physics problems often require interpreting textual instructions, identifying relevant principles, performing calculations, and evaluating results \autocite{Kontopoulou2024, Sweller2011}. Students with learning disabilities may find the cognitive demands of these multi-step tasks overwhelming, particularly when differences in executive functioning or information processing are present. Sweller \autocite{Sweller2011} emphasizes that managing high cognitive load is critical for student learning, and without appropriate scaffolding, students may experience frustration, decreased confidence, and disengagement.

In response to these challenges, assistive technologies and accessibility tools have been developed to support disabled students in STEM fields. Screen readers and text-to-speech software enable students with visual or reading-related disabilities to access digital materials auditorily \autocite{Wattanakasiwich2025, Alenezi2025}. MathML-compatible tools allow students to navigate mathematical notation by translating equations into structured verbal descriptions, enhancing comprehension of symbolic content \autocite{Wattanakasiwich2025, Alenezi2025}. Lannan, Scanlon, and Chini \autocite{Lannan2019} argue that integrating these tools directly into course design promotes accessibility while maintaining rigorous engagement with physics content.

Interactive simulations and visualization tools further enhance accessibility. Complex physical phenomena are often difficult to convey textually, but simulation platforms allow students to manipulate variables and observe system dynamics in real time, supporting learners who benefit from visual or exploratory approaches \autocite{Wattanakasiwich2025, Alenezi2025, Lannan2019}. Adaptive laboratory equipment, remote experimentation platforms, and digital data collection systems also enable students with physical or sensory disabilities to participate fully in experimental activities. Ref.~\cite{Lannan2019} argues that embedding these tools directly into course design promotes accessibility while maintaining rigorous engagement with physics content, and Refs.~\cite{Wattanakasiwich2025, Alenezi2025} document broader improvements in student engagement when multiple modalities are available --- benefits that extend to the general student population, not only those with identified disabilities.

These strategies align with Universal Design for Learning (UDL) principles, which advocate providing multiple means of representation, engagement, and expression to support diverse learners. Schreffler et al.\ \autocite{Schreffler2019} and Sakowicz and Hamidi \autocite{Sakowicz2025} note that incorporating UDL principles into STEM course design improves accessibility while preserving disciplinary rigor.

Finally, systemic barriers within educational institutions can limit the effectiveness of accessibility supports. Gull, Kaur, and Basha \autocite{Gull2025} and Almufareh et al.\ \autocite{Almufareh2025} highlight that minimal training in inclusive pedagogy, particularly for physics faculty, can result in instructional practices that unintentionally exclude students with diverse learning needs. Institutional policies may also fail to address challenges unique to physics, such as laboratory accessibility, specialized software, and mathematical notation \autocite{Gull2025, Almufareh2025}.

Together, these structural, cognitive, and social factors illustrate the complex educational landscape faced by students with disabilities in physics programs. Difficulties related to reading comprehension, mathematical processing, cognitive load, and social stigma interact with curriculum structure to create barriers that extend beyond individual differences. Understanding these challenges is essential for designing inclusive teaching practices and educational technologies that promote equitable participation in STEM disciplines.


\section{How AI is Being Utilized in Physics Education}
\label{sec:AI_STEMEd}
AI technologies are increasingly being integrated into physics education. These tools have the potential to improve accessibility, support diverse learning needs, and assist instructors in developing more flexible teaching materials. While much of the public discussion surrounding AI in education focuses on concerns about academic integrity, research literature increasingly highlights the ways in which AI can support both instructors and students. In particular, AI systems can generate visual explanations, support personalized learning environments, and assist with the development of course materials. These capabilities are especially beneficial for students with learning disabilities, while also improving learning outcomes for the broader student population \autocite{Voultsiou2025, HarkinsBrown2025, USDOE2023}.


\subsection{Faculty Use}

One of the most immediate applications of AI in physics education is its use by instructors to develop and organize teaching materials. Generative AI systems can assist faculty members in producing lecture slides, problem sets, worked examples, and conceptual explanations. In the study by Wattanakasiwich, Kaewkhong, and Katwibun \autocite{Wattanakasiwich2025}, physics instructors reported that generative AI tools allowed them to quickly generate multiple explanations of the same concept, enabling more tailored instructional materials for students with diverse backgrounds and varying levels of preparation. Similarly, Skogvoll and Odden \autocite{Skogvoll2025} highlight that these systems can serve as productivity aids, reducing the time required to prepare teaching resources while also supporting differentiated instruction. Faculty can leverage AI to create alternative representations of complex ideas, which research suggests may enhance conceptual understanding in diverse classrooms \autocite{Wattanakasiwich2025, Skogvoll2025}.

AI can also support instructors in generating visualizations and simulations that make abstract scientific concepts more accessible. Many physics topics involve phenomena that are challenging to convey through text alone. Ben-Zion et al.\ \autocite{BenZion2025} describe AI-powered tools capable of producing interactive simulations, diagrams, and visual models that illustrate concepts such as orbital motion, electromagnetic fields, or molecular interactions. In the context of special education, Voultsiou and Moussiades \autocite{Voultsiou2025} and Harkins-Brown, Carling, and Peloff \autocite{HarkinsBrown2025} demonstrate that multimodal learning environments---combining visual, auditory, and interactive content---improve comprehension by presenting information through multiple sensory channels. These technologies allow instructors to deliver complex material in formats that are both engaging and accessible.

AI applications also extend to course administration and assessment. Instructors can employ AI systems to provide automated feedback on assignments, organize course content, and identify areas where students struggle. According to Wattanakasiwich, Kaewkhong, and Katwibun \autocite{Wattanakasiwich2025}, automated feedback can handle routine grading tasks, allowing faculty to focus on higher-order guidance and conceptual mentoring. In large STEM courses, where instructors may oversee hundreds of students, these tools can significantly reduce administrative burden while enhancing instructional efficiency.

Another important application of AI involves curriculum design and instructional planning. Liang et al.\ \autocite{Liang2026} note that AI systems analyzing student performance data can identify concepts that consistently challenge learners, guiding course redesign efforts and suggesting where additional scaffolding or alternative explanations may be warranted. By supporting evidence-based adjustments, AI can help instructors improve the effectiveness of physics teaching practices without supplanting the essential mentorship and interactive problem-solving that human instructors provide.

Despite these advantages, research emphasizes that AI should primarily function as an assistive tool rather than a replacement for educators. The U.S. Department of Education \autocite{USDOE2023} underscores that effective STEM education relies on human guidance, collaborative discussion, and conceptual mentoring. AI is most beneficial when it enhances instructional capabilities, enabling instructors to personalize learning experiences while maintaining the core educational role of the teacher \autocite{USDOE2023, Sakowicz2025}.


\subsection{Student Use}

Artificial intelligence technologies provide substantial benefits directly to students, particularly through personalized learning support and adaptive tutoring systems. AI-based tutoring platforms can guide students through complex problem-solving processes, provide hints, and offer immediate feedback on errors. These systems can decompose complicated problems into smaller, more manageable steps, helping students reduce cognitive load associated with multi-stage reasoning tasks common in STEM subjects \autocite{Paglialunga2025, Khowaja2025}. Adaptive learning systems can also adjust the pace and difficulty of instruction based on individual student performance, allowing learners to dedicate additional time to challenging concepts while progressing more efficiently through material they already understand. Research on AI-based adaptive learning technologies emphasizes that these systems can create inclusive learning environments by tailoring instruction to students' specific cognitive needs, thereby supporting students with disabilities who might otherwise struggle in standardized classroom settings \autocite{Paglialunga2025, Khowaja2025}.

AI tools also enhance students' ability to visualize abstract concepts through interactive simulations, diagrams, and dynamic models. Physics and other STEM disciplines frequently require learners to construct mental models of processes that cannot be directly observed, such as electric field interactions, orbital motion, or quantum phenomena \autocite{Finkelstein2005}. AI-generated simulations allow students to manipulate variables, observe outcomes in real time, and explore the relationships between different physical quantities. Studies of AI, virtual reality, and immersive technologies in special education suggest that multimodal representations improve conceptual understanding for students with learning disabilities who may struggle with text-based explanations alone \autocite{Voultsiou2025, HarkinsBrown2025}. For example, in quantum mechanics, interactive visualizations of a particle in a potential well can demonstrate how the wavefunction evolves with changes in boundary conditions or potential depth, supporting comprehension in ways that static equations cannot \autocite{Finkelstein2005}.

Research in physics education has consistently shown that visual and interactive learning tools improve conceptual understanding in complex subjects such as quantum mechanics and electromagnetism. These tools help students build more accurate mental models of abstract phenomena by connecting mathematical formalism with dynamic visual representations \autocite{Finkelstein2005}. AI-generated visualizations provide an additional layer of accessibility, particularly for learners with cognitive differences that make processing dense text or symbolic notation challenging.

Another significant advantage of AI systems is that they offer academic support in private and nonjudgmental environments. Students with learning disabilities often face stigma or hesitation when requesting help in public classroom settings or during office hours. AI-based tutoring systems allow students to ask questions, request clarifications, and receive explanations without disclosing disabilities to instructors or peers \autocite{Khowaja2025}. This anonymity can reduce psychological barriers to seeking support, increase engagement, and promote active learning.

AI tools also assist students in reading comprehension and academic writing, which are sometimes underemphasized in STEM education. Because physics curricula often prioritize quantitative problem solving, students may receive limited training in writing, summarizing, and explaining technical ideas \autocite{Alenezi2025}. Large language models can support learners by summarizing dense readings, clarifying unfamiliar terminology, and providing feedback on written work, including laboratory reports, research summaries, or conceptual explanations. Comparative studies of AI use in physics versus humanities suggest that AI writing tools may bridge the gap between technical knowledge and effective written communication, helping students articulate complex ideas more clearly \autocite{Alenezi2025}.

AI tutoring systems further reduce cognitive load by guiding students through the multiple stages of problem solving. Physics problems frequently require simultaneous processing of textual descriptions, mathematical expressions, and conceptual reasoning. AI platforms can support learners by highlighting relevant principles, suggesting strategies for constructing mathematical models, and offering stepwise interpretation of results \autocite{Paglialunga2025, Sweller2011}. This scaffolding allows students to maintain focus, build confidence, and develop independent problem-solving skills without becoming overwhelmed by the complexity of tasks.

An additional consideration highlighted in research is the sometimes ``adversarial'' dynamic between physics instructors and students. Traditional physics instruction often emphasizes rapid problem-solving, adherence to formal procedures, and strict evaluation of conceptual correctness, which can create tension for learners who struggle with conventional approaches. AI-supported learning environments can mitigate this tension by providing individualized guidance and alternative representations that accommodate diverse cognitive strategies, thereby reducing the adversarial perception of physics instruction and allowing students to engage more productively with course material \autocite{Paglialunga2025, Khowaja2025}.

In conclusion, AI technologies offer students direct benefits by providing adaptive, personalized support, reducing cognitive load, enhancing visualization and conceptual understanding, and allowing private, low-stakes access to learning resources. These tools also help students navigate challenging communication and writing tasks and can mediate the adversarial dynamics often present in physics instruction. When implemented thoughtfully, AI systems complement traditional teaching methods by promoting inclusive, accessible, and student-centered learning experiences \autocite{Paglialunga2025, Voultsiou2025, HarkinsBrown2025, Sweller2011, Khowaja2025, Alenezi2025, Finkelstein2005}.


\section{How AI Can Support Students with Disabilities in Physics Education}
\label{sec:AI_Support}
Artificial intelligence tools offer several distinct advantages for students with disabilities beyond the general benefits they provide to all learners. In particular, AI technologies can facilitate experimentation with learning strategies, provide iterative explanations, assist with reading and writing tasks, reduce cognitive load, and enhance conceptual understanding through dynamic visualizations. These capabilities contribute to learning environments that are flexible, adaptive, and accessible for students with diverse cognitive profiles. Voultsiou and Moussiades \autocite{Voultsiou2025}, Harkins-Brown et al.\ \autocite{HarkinsBrown2025}, and the U.S. Department of Education \autocite{USDOE2023} note that AI systems in education can support personalized instruction while improving accessibility for a wide range of learners.

A common critique of AI systems is that they often function as ``black boxes,'' meaning their internal processes and decision-making pathways are not easily interpretable by users \autocite{Skogvoll2025, BenZion2025}. In traditional educational contexts, this lack of transparency is often seen as a limitation because it may obscure the reasoning behind a given solution or explanation \autocite{Skogvoll2025}. However, this perspective overlooks an important parallel: many students with learning disabilities experience their own cognitive processes as opaque or difficult to articulate. For example, a student with dyslexia or ADHD may struggle to explain how they understand a concept, how they approach a problem, or why certain strategies work for them \autocite{Almufareh2025, HarkinsBrown2025}. In this sense, the student's mind functions like a ``black box'' in which inputs (instructions, readings, or problems) are transformed into outputs (solutions or conceptual understandings) in ways that are not immediately visible to instructors. AI systems, when used thoughtfully, can interact with students' cognitive ``black boxes'' by providing iterative feedback, multiple representations of concepts, and flexible problem-solving pathways, thereby reducing barriers associated with verbalizing learning needs \autocite{Voultsiou2025, Khowaja2025}. This reframing highlights that a ``black box'' is not inherently negative; rather, AI can act as a scaffold that helps students explore and reveal their internal cognitive processes on their own terms.

For students with learning disabilities, articulating how they understand or struggle with a topic can create barriers to receiving effective support \autocite{Almufareh2025, HarkinsBrown2025}. AI-based learning systems can mitigate this by allowing learners to experiment with explanations, examples, and problem-solving strategies without needing to immediately justify their reasoning to another person \autocite{Khowaja2025, Weijers2025}. Research indicates that AI-supported environments promote iterative exploration, enabling students to test different approaches until they identify methods that align with their cognitive needs \autocite{Voultsiou2025, Khowaja2025}. In doing so, students develop self-awareness about how they best process information and normalize diverse learning strategies \autocite{Voultsiou2025, Weijers2025}. Adaptive AI systems can analyze patterns in student interactions and adjust feedback or instructional strategies accordingly, supporting personalized learning while reducing cognitive stress \autocite{Khowaja2025, Skogvoll2025}.

Another advantage of AI-supported learning systems is their continuous availability. Unlike office hours or tutoring sessions, which occur at fixed times, AI systems allow students to revisit explanations, request additional examples, and explore concepts at their own pace \autocite{Khowaja2025}. This flexibility is particularly valuable for students who require additional processing time to comprehend complex material or revisit difficult concepts multiple times before achieving full understanding. AI-powered summarization and explanation tools can translate dense technical readings into more accessible language, clarify unfamiliar terminology, and scaffold comprehension for students who struggle with language processing \autocite{Voultsiou2025, HarkinsBrown2025}. AI writing assistants similarly support students in organizing written work, refining scientific explanations, and iteratively revising reports, summaries, and conceptual descriptions while receiving structured feedback \autocite{Voultsiou2025, HarkinsBrown2025}.

AI technologies are uniquely capable of supporting diverse cognitive profiles because they can dynamically adapt instruction to individual learning preferences. Unlike traditional assistive technologies that target specific impairments, AI systems can provide customized explanations, pacing adjustments, and alternative representations of the same material. For instance, a single physics concept may be presented through textual explanations, visual diagrams, mathematical formulations, interactive simulations, or step-by-step problem-solving guides \autocite{Voultsiou2025, Khowaja2025}. Al-Azawei \autocite{AlAzawei2023} notes that such adaptive technologies can significantly enhance participation and engagement for students with disabilities when implemented effectively.

A central benefit of AI-supported learning is its ability to reduce cognitive load during complex tasks. STEM subjects often require simultaneous processing of multiple layers of information, including conceptual explanations, mathematical formulas, and graphical representations. AI systems can scaffold learning by breaking problems into smaller, manageable steps, providing targeted hints, and guiding students through iterative problem-solving sequences. This scaffolding allows learners to maintain focus, reduce cognitive overload, and build confidence when engaging with challenging material. Paglialunga and Melogno \autocite{Paglialunga2025} and Sweller \autocite{Sweller2011} emphasize that AI interventions are most effective when they support active problem solving rather than replace it entirely, balancing guidance with intellectual challenge to foster meaningful learning.

Visualization technologies powered by AI can further enhance conceptual understanding, particularly in abstract domains of physics \autocite{Liang2026}. Many scientific concepts, such as fields, waves, or quantum states, require students to construct mental models of processes that cannot be directly observed. AI-generated simulations, interactive diagrams, and dynamic visualizations allow learners to explore these concepts from multiple perspectives and observe the effects of varying parameters in real time \autocite{Voultsiou2025, HarkinsBrown2025, Liang2026}. These multimodal approaches improve engagement and help students form more accurate mental models, particularly for those who benefit from spatial or visual learning strategies \autocite{Voultsiou2025, HarkinsBrown2025}.

A concrete illustration of these benefits can be found in quantum mechanics, a field notoriously challenging for students due to its abstract nature. Understanding the particle in a potential well, for example, requires students to interpret the Schr\"{o}dinger equation, apply boundary conditions, and visualize standing wave solutions \autocite{Singh2001}. For students with learning disabilities, this combination of mathematical formalism and conceptual abstraction can be particularly difficult. AI-powered systems can present the same concept in multiple ways: generating dynamic visualizations of the wavefunction as potential parameters change, providing analogies to vibrating strings or standing waves, or offering step-by-step derivations of relevant equations \autocite{Singh2001, Finkelstein2005}. By iteratively exploring these representations, students can identify explanations that best align with their cognitive strengths, improving conceptual understanding without the pressure of immediately verbalizing their reasoning in traditional classroom settings. In this sense, AI tools act as both a bridge to abstract concepts and a personalized cognitive scaffold for learners with diverse needs.


\section{Concerns}
\label{sec:concerns}
Although artificial intelligence technologies offer numerous opportunities for improving accessibility and supporting diverse learning needs, their integration into physics education also raises several important concerns. Researchers emphasize that AI tools must be implemented carefully to ensure they support learning rather than replace critical reasoning processes. While AI systems can assist with explanations, visualizations, and adaptive tutoring, they also introduce potential risks related to overreliance, misinformation, bias, and technical limitations. These issues are particularly salient for students with disabilities, who may interact with AI tools in ways that differ from the general student population and whose learning needs require careful scaffolding \autocite{Paglialunga2025, Voultsiou2025}.

One primary concern involves the use of generative AI to complete physics assignments and problem sets. Physics problem solving requires multi-step reasoning, careful interpretation of mathematical relationships, and the ability to evaluate whether solutions are physically meaningful. Current generative AI systems often struggle with these tasks, particularly when problems involve complex derivations, iterative reasoning, or multiple layers of abstraction. For instance, AI-generated solutions to quantum mechanics problems, such as those involving the Schr\"{o}dinger equation, may produce mathematically plausible results that contain subtle conceptual errors, misapplied boundary conditions, or incorrect interpretations of wavefunction behavior \autocite{Wattanakasiwich2025, Weijers2025, Singh2001}. Several studies therefore emphasize that AI should not be used to directly complete assignments; instead, it is most effective when employed to provide hints, stepwise guidance, or conceptual explanations that support active problem solving \autocite{Wattanakasiwich2025, Weijers2025}.

Another major concern is the potential for excessive reliance on AI tools. AI systems can provide rapid answers, adaptive explanations, and simplified summaries, which may inadvertently reduce the incentive for students to engage deeply with complex material \autocite{Paglialunga2025, Voultsiou2025}. This issue can be particularly pronounced for students with learning disabilities, who may rely on AI for accessible explanations or alternative learning strategies. For example, a student struggling to visualize the time evolution of a particle in a quantum potential well may repeatedly request AI-generated simulations rather than attempt intermediate analytic steps or mental modeling. While these visualizations can aid understanding, overdependence may limit opportunities for students to practice independent reasoning, evaluate conceptual consistency, or develop problem-solving resilience. In this context, AI functions as both a cognitive scaffold and a potential crutch; careful guidance from instructors is needed to balance support with active engagement \autocite{Paglialunga2025, Voultsiou2025, Khowaja2025, Singh2001}.

Students with disabilities face additional challenges related to AI reliance. Because they may experience their own cognitive processes as opaque or difficult to articulate, there is a risk that AI explanations could be accepted uncritically if they appear to match the student's immediate understanding \autocite{Voultsiou2025, Khowaja2025}. For example, an AI-generated visualization of a wavefunction may help a student quickly see qualitative features of the system, but if it misrepresents boundary conditions or node structures, the student may internalize incorrect patterns. Without deliberate scaffolding, these errors can compound, particularly for learners who struggle with abstract mathematical reasoning, working memory, or executive function. This highlights the need for AI to be integrated within structured learning environments that encourage verification, reflection, and iterative practice \autocite{Paglialunga2025, Voultsiou2025, Weijers2025}.

Automated AI-detection systems present another concern. Many universities have implemented software to identify AI-generated writing, but research shows that these systems frequently produce false positives, labeling human-written work as AI-generated \autocite{Voultsiou2025, HarkinsBrown2025}. This issue may disproportionately affect students with disabilities, particularly those who rely on assistive writing technologies, text-to-speech systems, or structured composition tools. Differences in style, syntax, or phrasing introduced by accessibility tools can trigger detection algorithms, potentially leading to unfair academic penalties. Scholars caution that reliance on automated detection without human oversight could exacerbate existing inequities for students who already face accessibility barriers \autocite{Voultsiou2025, HarkinsBrown2025}.

Generative AI systems are also susceptible to bias and hallucinations, producing inaccurate explanations, fabricated references, or misleading content \autocite{Paglialunga2025, Voultsiou2025, Weijers2025}. While these errors affect all users, they pose additional risks for students with disabilities who may place greater trust in AI-generated explanations due to the accessibility and clarity of the outputs. For example, a student with dyslexia or ADHD might prefer AI-generated textual summaries of complex physics derivations but may not have the metacognitive resources to detect subtle errors in reasoning. This underscores the importance of teaching students to critically evaluate AI outputs and to use these tools as complementary supports rather than authoritative sources \autocite{Paglialunga2025, Voultsiou2025, Weijers2025}.

Technical limitations of AI accessibility tools present further challenges. Many systems rely on video analysis, speech recognition, or gesture-based interaction, which may perform poorly for students whose behaviors or communication patterns differ from those represented in training datasets. Neurodivergent learners or students with motor differences may encounter difficulties using AI tools as intended, reducing reliability and potentially excluding the very populations these technologies aim to assist \autocite{Voultsiou2025, HarkinsBrown2025}.

Finally, broader societal and institutional factors influence AI integration in physics education. Public discussions of AI often emphasize risks such as automation, misinformation, and academic dishonesty, which can create hesitation among instructors and institutions. Overly restrictive policies may prevent faculty from exploring AI applications that could improve accessibility and learning outcomes \autocite{Liang2026}. Additionally, many STEM instructors receive limited training in educational technology or inclusive pedagogy, limiting their ability to effectively guide students in responsible AI use \autocite{Almufareh2025, Voultsiou2025}. Successful integration therefore requires coordinated efforts among instructors, institutions, and students, with clear guidelines for equitable and responsible AI adoption.

In summary, AI technologies offer powerful opportunities to support learning, but they introduce new challenges, particularly for students with disabilities. Risks such as overreliance, cognitive misalignment, accessibility limitations, and biased outputs require careful mitigation. Concrete examples, such as AI-assisted visualization of the Schr\"{o}dinger equation, illustrate both the potential benefits and hazards: while AI can provide accessible, dynamic representations of wavefunctions, students must still engage in active reasoning to verify conceptual and mathematical correctness. When implemented thoughtfully, AI can enhance physics education and accessibility, but without structured guidance, these tools may inadvertently reinforce misconceptions or inequities \autocite{Paglialunga2025, Voultsiou2025, Khowaja2025, Weijers2025, Singh2001}.


\section{Discussion and Summary}
\label{sec:discussion}
This study examined the role of artificial intelligence (AI) within undergraduate physics education, focusing on its implications for accessibility, learning, and equitable participation. The analysis was guided by three research questions from the original proposal: (1) how generative AI is currently utilized in undergraduate physics education, (2) how these technologies are perceived by students and faculty, and (3) what best practices can support the responsible integration of AI tools within physics learning environments.

The findings suggest that AI technologies are already influencing multiple aspects of physics education, but their impact depends heavily on how they are implemented and supported by both instructors and institutions.


\subsection{RQ1: Utilization of AI in Undergraduate Physics Education}

The first research question explored how generative AI is being used in undergraduate physics courses. Evidence indicates that students primarily employ AI tools as conceptual assistants rather than as automated problem solvers. Students often use AI to translate between different representations of physics problems, such as converting textual descriptions into equations, generating diagrams, or clarifying derivations. These uses reflect the convergent and rule-based nature of physics problem solving, in which solutions must conform to established physical laws rather than subjective interpretation.

Differences also appear between lower-year and upper-year courses. In introductory courses, AI is often used as a supplementary tutor, providing step-by-step explanations for mechanics, kinematics, or electromagnetism problems. Here, AI can scaffold reasoning without replacing instructor guidance, helping students build foundational skills in interpreting physical descriptions and performing calculations. In upper-year courses, particularly those involving computational physics, statistical mechanics, or quantum mechanics, AI is increasingly employed to assist with symbolic manipulation, coding tasks, and interpreting numerical simulations.

This progression highlights a potential solution: integrating AI instructionally across the curriculum rather than leaving its use unstructured. For lower-year students, AI tools should emphasize guided exploration and conceptual scaffolding, while upper-year students may benefit from AI-assisted simulation, modeling, and advanced problem-solving exercises. Embedding AI-based activities into course design, such as asking students to critically evaluate AI-generated derivations or compare multiple AI-provided solutions, can help maintain analytical rigor while enhancing accessibility and comprehension for all learners, including those with disabilities.

\subsection{RQ2: Perceptions of AI Among Students and Faculty}

Student perceptions of AI in physics education are generally positive, particularly for personalized, on-demand support. Students report that AI assists with conceptual understanding, reduces cognitive load, and allows exploration of alternative problem-solving strategies at their own pace. For students with learning disabilities, these benefits can be especially meaningful, providing accessible explanations, reducing barriers associated with dense technical readings, and allowing private engagement with material that might otherwise be intimidating in classroom settings.

Faculty perspectives tend to be more cautious. Instructors frequently cite concerns about academic integrity, overreliance, and subtle AI-generated errors in conceptual or mathematical reasoning. Physics education is particularly sensitive to small mistakes, where a single incorrect assumption can render an entire solution invalid. These concerns emphasize the importance of structured integration, where AI serves as a learning facilitator rather than a replacement for critical thinking.

Successful implementation depends heavily on instructors' familiarity with AI tools and their ability to guide students in responsible use. Faculty who understand how AI works can design assignments that integrate AI thoughtfully, such as using AI-generated explanations for peer review, error analysis, or conceptual comparisons. Professional development for instructors is therefore critical, as it enables them to scaffold AI interactions, anticipate common student misunderstandings, and ensure that tools support rather than undermine independent reasoning. This aligns directly with later recommendations regarding instructor training, ethical guidance, and accessibility evaluation.


\subsection{RQ3: Best Practices for Integrating AI into Physics Education}

The literature suggests several emerging best practices for AI integration:

\textbf{Guided AI Support:} AI tools are most effective when they provide hints, explanations, or visualizations that supplement rather than replace problem solving. This approach allows students to retain responsibility for analytical reasoning while benefiting from scaffolding and feedback.

\textbf{Curricular Integration:} AI should be explicitly embedded in course activities. For example, students may be asked to identify errors in AI solutions, compare multiple approaches, or generate AI-assisted visualizations for conceptual understanding.

\textbf{AI Literacy for Students and Instructors:} Both groups need training in evaluating AI outputs to avoid misconceptions and overreliance. Students benefit from understanding how AI generates answers, while instructors gain confidence in designing assignments that leverage AI responsibly.

\textbf{Accessibility Considerations:} AI systems should be evaluated for inclusivity, ensuring that students with diverse cognitive and communication needs can engage effectively. This includes testing for biases in explanations, reading support, and visualizations.

These practices foreshadow specific recommendations later in the paper, highlighting the interplay between technology, pedagogy, and accessibility. They suggest that AI's potential is greatest when it is structured, monitored, and paired with instructional guidance, particularly for students with learning disabilities who may rely on adaptive supports more heavily.


\subsection{Implications for Accessible Physics Education}

The findings from these research questions suggest that artificial intelligence has the potential to play a meaningful role in improving accessibility within physics education. When used appropriately, AI tools can provide flexible learning environments that support diverse cognitive styles, offer multiple representations of complex concepts, and reduce barriers associated with technical reading and multi-step problem solving. These features may be particularly beneficial for students with learning disabilities who experience difficulties with traditional instructional formats.

However, the literature also emphasizes that the integration of AI technologies must be approached carefully. Without appropriate guidance, AI systems may reinforce misconceptions, encourage passive learning, or introduce new forms of inequity related to technological access and digital literacy. Effective implementation therefore requires a balanced approach that combines technological innovation with inclusive teaching practices.


\section{Recommendations}
\label{sec:recommendations}
\begin{enumerate}
    \item \textbf{Integrate AI as a Learning Support Tool Rather Than a Solution Generator:}
    AI technologies should scaffold problem-solving and conceptual understanding without producing complete solutions. Students should be encouraged to critically engage with AI outputs while maintaining responsibility for their reasoning and calculations.

    \item \textbf{Provide Training for Instructors on AI-Supported and Inclusive Pedagogy:}
    Professional development programs should equip instructors with AI literacy and guidance on supporting diverse learning needs. Faculty who understand AI can design assignments and activities that enhance learning while preventing overreliance or passive consumption.

    \item \textbf{Establish Institutional Guidelines for Ethical and Accessible AI Use:}
    Universities should define clear policies for legitimate AI use, including distinguishing between learning support and academic misconduct. Policies should also protect students who use assistive technologies from bias in automated AI-detection systems.

    \item \textbf{Evaluate AI Tools for Accessibility and Bias:}
    Before implementing AI-based educational technologies, institutions should test systems with diverse users to ensure they are accessible and equitable. Bias in AI explanations, visualizations, or interaction modes should be addressed proactively.

    \item \textbf{Encourage Research on AI-Supported Learning in Physics Education:}
    Further studies should examine how AI impacts learning outcomes, particularly for students with disabilities. Collaborative research involving physics educators, special education experts, and educational technologists can help identify best practices for inclusive AI integration.
\end{enumerate}

\printbibliography

\appendix

\section{Appendices and Resources}

\end{document}